\section{Предложенный метод}
\label{sec:proposed_method}

В данном разделе представлено описание разработанного метода адаптивной растеризации поверхностей, заданных знаковыми функциями расстояния (SDF), с использованием иерархических структур данных SComTree. Предложенный подход базируется на гибридной схеме, сочетающей эффективность разреженных структур данных и производительность современного графического конвейера. Основной алгоритм разделен на четыре ключевых этапа: формирование массива активных узлов, триангуляция поверхности, растеризация и постобработка.

\subsection{Формирование массива активных узлов}

Первым и наиболее критическим этапом является идентификация тех фрагментов пространства, которые содержат участки искомой поверхности $\Omega$ и попадают в область видимости. В основе этого этапа лежит иерархический обход октодерева (или DAG в случае SComTree), дополненный предикатами отсечения.

Алгоритм обхода начинается от корня структуры и спускается к листьям, выполняя для каждого узла следующие проверки:
\begin{enumerate}
    \item \textbf{Геометрическое отсечение (Frustum Culling):} Узел отбрасывается, если его ограничивающий объем (AABB) полностью находится за пределами пирамиды видимости.
    \item \textbf{Окклюзионное отсечение (Occlusion Culling):} С использованием иерархического буфера глубины (Hi-Z) проверяется, не скрыт ли узел за уже отрисованными объектами. Это позволяет избежать детальной обработки геометрии, которая не будет видна на итоговом изображении.
    \item \textbf{Анализ пересечения с поверхностью:} Узел считается потенциально содержащим поверхность, если диапазон значений SDF внутри него включает ноль. В формате SComTree это определяется на основе типов дочерних узлов и их связей.
\end{enumerate}

Результатом данного этапа является линейный массив \textit{активных узлов} (Active Nodes/Leafs), каждый из которых представляет собой кубическую область пространства определенного уровня детализации. Использование линейного массива на выходе позволяет эффективно распределять нагрузку на последующих стадиях.

\subsection{Триангуляция поверхности}

После того как список активных узлов сформирован, выполняется локальное извлечение поверхности внутри каждого такого узла. Для этого используется алгоритм Marching Cubes (или его модификация Marching Tetrahedra для избежания топологических неоднозначностей).

Процесс триангуляции для каждого узла включает:
\begin{itemize}
    \item \textbf{Выборку значений SDF:} В восьми вершинах кубического узла вычисляются (или извлекаются из структуры данных) значения знаковой функции расстояния.
    \item \textbf{Интерполяцию позиций вершин:} Если значения SDF в соседних вершинах имеют разные знаки, это означает, что поверхность пересекает ребро между ними. Точное положение вершины треугольника $P$ на ребре $(P_1, P_2)$ вычисляется с помощью линейной интерполяции:
    \begin{equation}
        P = P_1 + (P_2 - P_1) \cdot \frac{0 - SDF(P_1)}{SDF(P_2) - SDF(P_1)}
    \end{equation}
    \item \textbf{Генерацию нормалей:} В каждой извлеченной вершине вычисляется вектор нормали к поверхности как нормированный градиент SDF. Применение метода центральных разностей позволяет получить гладкие нормали даже для дискретных представлений.
\end{itemize}

Итогом этапа является набор треугольников, локально аппроксимирующих неявную поверхность с заданной точностью.

\subsection{Растеризация}

Сгенерированные на предыдущем этапе треугольники подаются на вход стандартного этапа растеризации. Однако, в отличие от классического рендеринга статических сеток, здесь применяется концепция \textit{отложенного рендеринга} (Deferred Rendering) и использование G-буфера (Geometry Buffer).

Растеризатор преобразует трехмерные треугольники в двумерные фрагменты (пиксели), заполняя G-буфер следующими данными:
\begin{itemize}
    \item \textbf{Мировые координаты или глубина:} Для последующего восстановления положения точки в пространстве.
    \item \textbf{Вектор нормали:} В экранном пространстве или мировых координатах для расчета освещения.
    \item \textbf{Материал и цвет:} Атрибуты, специфичные для конкретного объекта.
\end{itemize}

Использование G-буфера позволяет отделить геометрическую сложность сцены от сложности освещения, что особенно важно при работе с процедурно генерируемой геометрией.

\subsection{Постобработка и освещение}

На финальном этапе выполняется расчет итогового цвета каждого пикселя на основе данных из G-буфера. Этот этап включает:
\begin{enumerate}
    \item \textbf{Расчет освещения (Shading):} Применение моделей освещения (например, модели Фонга или PBR) с использованием нормалей и позиций из G-буфера.
    \item \textbf{Глобальное затенение (Ambient Occlusion):} Для SDF-представлений расчет Ambient Occlusion может быть выполнен более эффективно за счет информации о расстоянии до ближайших поверхностей, заложенной в самой структуре SComTree.
    \item \textbf{Сглаживание и коррекция:} Применение фильтров постобработки (FXAA, TAA) для устранения ступенчатых артефактов, характерных для воксельных и иерархических структур.
\end{enumerate}

Таким образом, предложенный метод обеспечивает сквозной цикл визуализации неявных поверхностей, эффективно масштабирующийся от грубых приближений до высокодетализированных моделей.